% ---------------------------------------------------------------
% INTRODUCTION
% ---------------------------------------------------------------
\chapter*{Úvod}\addcontentsline{toc}{chapter}{Úvod}\markboth{Úvod}{Úvod}
\setcounter{page}{1}

Přestože o nich většina uživatelů internetu nemá ponětí, datová centra tvoří nepostradatelný základ moderního digitálního světa. Jejich role se postupně vyvinula od jednoduchých úložišť dat až po komplexní výpočetní uzly v měřítku celých skladových hal. Zlomovým okamžikem v jejich evoluci se stal nástup generativní umělé inteligence (AI), který radikálně změnil nároky nejen na samotná datová centra, ale také na energetickou infrastrukturu. Oproti tradičním cloudovým službám totiž trénování a provoz moderních AI modelů vyžaduje extrémní výpočetní výkon.

Současná AI infrastruktura čelí kritice kvůli svým environmentálním dopadům, které zahrnují nejen obrovskou spotřebu elektřiny, ale také významné emise skleníkových plynů a narůstající odběr vody. Odhaduje se, že vývoj, trénování a provoz velkých jazykových modelů může během jednoho roku používání vyprodukovat až 160 000 tun uhlíku (jako provoz 30 000 benzinových aut) \cite{jegham_how_2025} a odebrat asi 5 miliard litrů sladké vody (což je polovina odběru Spojeného království) \cite{li_making_2025}. Tato technologická náročnost vytváří bezprecedentní tlak na lokální zdroje v místech, kde se AI datová centra koncentrují.

Ohniskem rozvoje AI infrastruktury jsou v současnosti Spojené státy americké, kde se v nadcházejících letech očekává prudký nárůst poptávky po elektřině vyvolaný zvýšenou poptávkou po výpočetním výkonu. Tento trend se však postupně globalizuje a zasahuje i Evropu. Vzhledem k tomu, že se otevření prvního AI datového centra v České republice očekává již na konci roku 2026 a zprovoznění největšího tuzemského DC je plánováno na začátek roku 2027, je nezbytné včas analyzovat rizika a příležitosti s nimi spojené.

Tato práce slouží jako studie socioekonomických dopadů investic do AI infrastruktury na území USA, které v této oblasti představují nejrozvinutější trh. Cílem je na základě analýzy amerických dat identifikovat klíčové korelace a vytvořit prediktivní model, který umožní simulovat dopady výstavby podobných center v České republice a potenciálně i v dalších zemích čelících AI transformaci. Cílem naopak není hodnotit přínosy umělé inteligence jako takové -- model má sloužit k prozkoumání dopadů \textit{fyzické přítomnosti} datových center na území ČR, aniž by zahrnoval jejich produkty.

% Mozna by stalo zato rozdelit cile na teoretickou (definice, sebtani dat, analyza problematiky) a praktickou cast (vytvoreni modelu, vytvoreni aplikace). Dulezite upresnit co ma byt vystupem: webova stranka ktera vizualizuje data modelu
% Do prezentace vlozit taky otazky od oponenta, lze tam v odrazkach pridat odpoved